% paper_draft.tex — in-progress section drafts.
% Single source of truth for paper text once §5.5 / abstract / §6 unblock.
% Will be merged into final paper.tex at the end. Sections that need data
% are marked at top. Sections marked READY can be lifted as-is.
%
% Status (2026-06-01, post seq=51 LR probe):
%   §5.5 SGD pressure test:    DRAFTED below (3 subsections)
%   §6 "Why SGD fails":        DRAFTED below (mechanism speculation unlocked)
%   Abstract (§5.5 sentence):  DRAFTED below
%   Title:                     "SGD-Unreachable" preserved (see paper_diffs.md flip log)
%
% Everything else lives in user-written paper sections; this file is the
% delta against the original PAPER.md / paper.tex skeleton.

%==============================================================
% Abstract — §5.5 sentence replacement
%==============================================================
% Replace the existing "yet SGD-trained low-rank adapters reading that input
% recover essentially none of it ..." clause with:

\newcommand{\abstractsgdsentence}{%
yet SGD-trained low-rank adapters reading the \emph{same} input recover at
most $27\%$ of the residual at the most accessible layer and stay negative
at the deepest, across a $100\times$ learning-rate sweep
($10^{-5}$ to $3{\times}10^{-4}$) covering the standard LoRA default; and
even when initialized from the closed-form ridge solution, AdamW
finetuning at any tested learning rate degrades the solution monotonically,
indicating that the ridge optimum is not a local minimum of the
cross-entropy objective.%
}

%==============================================================
% §5.5 SGD pressure test (3 subsections, ~1.5 pages target)
%==============================================================

\subsection{Random-initialized SGD fails to find the ridge solution across a
100$\times$ learning-rate range}
\label{sec:sgd-random}

The closed-form ridge solution of \S\ref{sec:ridge} establishes that the
context residual is a generalizing linear function of $\mathrm{LN}_l(h)$. The
remaining question is whether gradient training of an adapter with the same
parameterization and the same input can reach the same recovery from random
initialization. If it could, the closed-form solution would be a useful
shortcut but not a methodological correction. If it cannot, the closed-form
solution is the deployable form and the gradient-trained adapter is a
diagnostic of optimization difficulty rather than of representational
content.

We sweep AdamW learning rate over $\{10^{-5},\,3{\times}10^{-5},\,10^{-4},\,
3{\times}10^{-4}\}$---a $100\times$ range that includes the LoRA-default
$10^{-4}$ \citep{hu2022lora}---at each of the three most context-bound
Pythia-410M layers ($L_6, L_7, L_{11}$), training for $5000$ steps at rank
$64$ and cap $0.5$ on top of the $A_1$ anchor (the same architecture and
budget under which the closed-form ridge attains
$R_\mathrm{context} \in [0.56, 0.78]$).
Table~\ref{tab:lr-probe-random} reports $R_\mathrm{context}$ on held data.

\begin{table}[t]
\centering\small
\caption{\textbf{Random-init AdamW fails to find the ridge solution across
$100\times$ learning rate}. $R_\mathrm{context}$ at rank $64$, cap $0.5$,
$5000$ steps; A1 anchor base. Closed-form ridge at the same rank and cap
(no SGD) gives $R_\mathrm{context} = 0.78, 0.74, 0.56$ at $L_6, L_7, L_{11}$.}
\label{tab:lr-probe-random}
\begin{tabular}{lcccc}
\toprule
layer & $\mathrm{lr}{=}10^{-5}$ & $3{\times}10^{-5}$ & $10^{-4}$ & $3{\times}10^{-4}$ \\
\midrule
$L_6$  &  0.27 & -0.11 &  0.05 & -0.19 \\
$L_7$  &  0.01 & -0.39 & -0.25 & -0.33 \\
$L_{11}$ & -1.05 & -1.17 & -1.26 & -1.13 \\
\bottomrule
\end{tabular}
\end{table}

No learning rate produces a random-initialized adapter that approaches the
closed-form ceiling. At $L_6$, the best random-init result is
$R_\mathrm{context} = 0.27$ at $\mathrm{lr}{=}10^{-5}$---roughly one-third
of the ridge's $0.78$. At $L_7$, the best is $0.01$. At the most
context-bound layer $L_{11}$, every random-init configuration is strongly
negative (worse than leaving the anchor untouched), with values in
$[-1.26,\,-1.05]$. The pattern is robust across optimizer and schedule
variations (Lion, linear-warmup, and curriculum norm-cap schedules tested
at $\mathrm{lr}{=}10^{-4}$ all fall within the envelope shown in
Fig.~\ref{fig:sgd-pressure}).

\paragraph{Why this is the right comparison.} The ridge solution and the
random-initialized adapter share \emph{the same input}
($\mathrm{LN}_l(h_{l-1})$), \emph{the same architectural class} (a
$d{\times}d$ map factored to rank $r{=}64$), and \emph{the same forward-time
norm budget} (cap $= 0.5$). They differ only in how the parameters are
chosen: by closed-form least squares with $\ell_2$ regularization, versus
gradient descent on next-token cross-entropy with $5{,}000$ AdamW steps. The
$\sim 0.7$ gap in $R_\mathrm{context}$ is therefore not a representation
limit, an expressivity bound, or a capacity issue; it is the gap that
gradient descent does not close.

\subsection{Ridge-initialized SGD also degrades, monotonically with learning rate}
\label{sec:sgd-ridge-init}

The natural next question is whether the failure in \S\ref{sec:sgd-random}
is an \emph{initialization} problem---perhaps gradient descent simply
cannot navigate to the ridge solution from a random start, but \emph{could}
maintain it from a good start. We test this directly by initializing the
LoRA factors $A, B$ to the SVD of the closed-form ridge solution and
running the same AdamW training at the same learning rate sweep.

Table~\ref{tab:lr-probe-ridge} reports the result. The data refutes the
initialization-only hypothesis: even starting from the closed-form optimum,
continued gradient training degrades the solution at every tested learning
rate above $10^{-5}$, and the degradation is monotonic with learning rate.
At $L_6$, ridge-init starts at $0.78$ and ends at $0.77$ when
$\mathrm{lr}{=}10^{-5}$ (essentially unchanged: $5{,}000$ steps at that
learning rate move the parameters negligibly), drops to $0.62$ at
$\mathrm{lr}{=}3{\times}10^{-5}$, $0.27$ at $10^{-4}$, and $-0.11$ at
$3{\times}10^{-4}$. The pattern is identical at $L_7$ ($0.73 \to 0.49 \to
-0.36 \to -0.27$) and $L_{11}$ ($0.44 \to -0.05 \to -1.20 \to -0.94$). At
all three layers, the larger the learning rate, the further the parameters
drift from the closed-form optimum and the worse the cross-entropy.

\begin{table}[t]
\centering\small
\caption{\textbf{Ridge-init AdamW also degrades the solution, monotonically
with learning rate}. Same setup as Table~\ref{tab:lr-probe-random} but
$A, B$ initialized from the closed-form ridge SVD. Top row is the
zero-shot closed-form recovery (no SGD) for reference.}
\label{tab:lr-probe-ridge}
\begin{tabular}{lcccc}
\toprule
 & $\mathrm{lr}{=}10^{-5}$ & $3{\times}10^{-5}$ & $10^{-4}$ & $3{\times}10^{-4}$ \\
\midrule
\textit{zero-shot}: $L_6{=}0.78$, $L_7{=}0.74$, $L_{11}{=}0.56$ & & & & \\
\midrule
$L_6$ ridge-init   &  0.77 &  0.62 &  0.27 & -0.11 \\
$L_7$ ridge-init   &  0.73 &  0.49 & -0.36 & -0.27 \\
$L_{11}$ ridge-init &  0.44 & -0.05 & -1.20 & -0.94 \\
\bottomrule
\end{tabular}
\end{table}

A subtle methodological observation reinforces the point. Our earlier
ridge-finetune measurement (\S\ref{sec:ridge}) reported ridge-init
$R_\mathrm{context} \approx 0.78/0.74/0.56$ at $L_6/L_7/L_{11}$ using
$500$ AdamW steps at $\mathrm{lr}{=}3{\times}10^{-5}$, suggesting the
solution was stable under SGD finetuning. The $5{,}000$-step probe at
$\mathrm{lr}{=}10^{-5}$ (covering the same effective $\mathrm{lr}\!\cdot\!
\mathrm{steps}$ product as the $500$-step measurement) yields
$0.77/0.73/0.44$. The $L_{11}$ discrepancy ($0.56 \to 0.44$, a $0.12$ drop)
is itself an instance of the slow drift: the closer the deeper layer's
ridge optimum lies to a non-minimum of the cross-entropy surface, the
more rapidly extended SGD walks away from it. Layers $L_6, L_7$ drift
slowly enough to look stable; $L_{11}$ does not. We do not interpret this
$0.12$ as noise.

\subsection{Mechanistic interpretation: the ridge optimum is not a local
minimum of the CE objective}
\label{sec:sgd-mechanism}

Both \S\ref{sec:sgd-random} and \S\ref{sec:sgd-ridge-init} are explained by
a single hypothesis: the closed-form ridge solution and a CE-loss local
minimum coincide only approximately, and the two objectives' gradient
landscapes are not aligned at the ridge point. Concretely:

The ridge solution
$W^\star = (X^\top X + \lambda I)^{-1} X^\top Y$ minimizes the
\emph{reconstruction} loss
$\mathbb{E}\,\|V^\mathrm{real} - (A_1 + X W)\|^2$
on the calibration distribution. By contrast, our SGD training minimizes
next-token cross-entropy with $W = AB$. The two are related but not
identical: CE penalizes value-projection errors that propagate to a
\emph{different} next-token distribution more heavily than errors that do
not, while reconstruction penalizes value-projection errors uniformly. If
the CE landscape's curvature at $W^\star$ is shallow or negative along
some directions in the rank-$r$ subspace, AdamW will descend those
directions until it finds a CE minimum that improves training loss but
hurts held-out generalization---and the descent is monotonic in the
effective step size.

This hypothesis predicts everything we observe:
\begin{itemize}[leftmargin=12pt,itemsep=0pt,topsep=2pt]
\item Ridge-init at very small lr (we test $10^{-5}$) preserves the ridge
solution: the SGD steps are too small to leave its basin within $5{,}000$
iterations, and the result matches the closed-form recovery on the
non-trough layers ($L_6, L_7$) and drifts only mildly at the trough
layer ($L_{11}$).
\item Ridge-init at larger lr leaks: AdamW finds the
shallow CE direction and drifts. The drift is monotonic in lr.
\item Random-init never finds the ridge basin at any lr: it descends to a
nearby CE minimum that is not the ridge optimum and is consistently worse
on held data, regardless of where on the lr range one starts.
\item Layer depth matters: $L_{11}$ has the highest token-anchor cost
($R_\mathrm{context}$ headroom is largest), and so has the most curvature
to lose---both ridge-init and random-init degrade most severely there.
\end{itemize}

We have not directly characterized the loss-surface geometry around the
ridge point---e.g., by measuring the Hessian conditioning of the CE
objective at $W^\star$ or by tracing the SGD trajectory in
$\|W - W^\star\|$ during training---and we flag this as the most natural
follow-up. What the present results establish is the
\emph{dissociation itself}: a generalizing linear readout of
$\mathrm{LN}_l(h)$ exists at every tested layer, is found by ordinary
ridge regression in closed form, and is not found or maintained by
gradient training from any initialization or learning rate we tested.

\begin{figure*}[t]
\centering
\includegraphics[width=\textwidth]{fig11.pdf}
\caption{\textbf{Across a $100\times$ learning-rate sweep, random-init
AdamW never reaches the closed-form ridge ceiling, and ridge-init AdamW
degrades monotonically with learning rate}. $R_\mathrm{context}$ at the
three most context-bound layers of Pythia-410M. Solid red: ridge-init
AdamW ($5{,}000$ step, $r{=}64$, cap${=}0.5$). Dashed blue: random-init
AdamW (same setup, fan-in init). Dotted red horizontal: closed-form ridge
solution evaluated without SGD (the deployable form). Markers (`$\times$'
random-init variants, `$+$' ridge-init) overlay the seq=49 six-variant
sweep (Lion, warmup, curriculum cap, 2k/5k-step AdamW) at
$\mathrm{lr}{=}10^{-4}$ ($3{\times}10^{-5}$ for Lion), confirming that
within-fixed-lr SGD variants stay inside the envelope.}
\label{fig:sgd-pressure}
\end{figure*}

%==============================================================
% §6 — first paragraph (mechanism speculation, now unblocked)
%==============================================================
% Replace the existing first paragraph of §6 Discussion with the version
% below. Subsequent paragraphs ("For interpretability methodology",
% "For value-path understanding", "As a method", "Why mid-stack") remain
% as in PAPER.md.

\subsection*{Discussion (\S6, replacement for current first paragraph)}

The central finding is a dissociation: for deep value contextualization,
the recoverable structure exists, is low-rank, and is linearly readable
from the layer-input hidden state---yet gradient descent neither finds
the ridge readout from random initialization nor maintains it from the
closed-form initialization. The latter is the sharper version: a
\emph{deployable, ground-truth} solution sits in parameter space, can be
written down, and is the optimum of the reconstruction objective that
defines a value correction---and AdamW on cross-entropy walks away from
it monotonically as the learning rate grows
(\S\ref{sec:sgd-ridge-init}, Table~\ref{tab:lr-probe-ridge}). The
conservative norm constraint adopted to stabilize SGD makes the gap worse
at the high-capacity regime but is not the source of the failure: the
soft cap is forward-time only, applies symmetrically to both
initializations, and is empirically dominated by the
\emph{learning-rate}-driven drift (Appendix~B). The optimization difficulty
is the bottleneck, and it manifests independently of how the parameters
are initialized.
